% Part: first-order-logic
% Chapter: natural-deduction
% Section: provability-consistency

\documentclass[../../../include/open-logic-section]{subfiles}

\begin{document}

\iftag{FOL}
      {\olfileid{fol}{ntd}{prv}}
      {\olfileid{pl}{ntd}{prv}}

\olsection{\usetoken{S}{derivability} आणि सुसंगतता}

आता आपण !!{derivability} संबंधाचे अनेक गुणधर्म सिद्ध करू. ते स्वतंत्रपणेही
महत्त्वाचे आहेत, पण संपूर्णता प्रमेयाच्या सिद्धतेत प्रत्येकाची भूमिका असेल.

\begin{prop}\ollabel{prop:provability-contr}
  जर $\Gamma \Proves !A$ आणि $\Gamma \cup \{!A\}$ विसंगत असेल, तर
  $\Gamma$ विसंगत आहे.
\end{prop}

\begin{proof}
$\Gamma$ पासून $!A$ ची !!{derivation}~$\delta_1$ आणि $\Gamma \cup \{!A\}$
पासून $\lfalse$ ची !!{derivation}~$\delta_2$ असू द्या. मग आपण पुढीलप्रमाणे
!!{derive} करू शकतो:
\begin{prooftree}
\AxiomC{$\Gamma, \Discharge{!A}{1}$}
\RightLabel{$\delta_2$}
\DeduceC{$\lfalse$}
\DischargeRule{\Intro{\lnot}}{1}
\UnaryInfC{$\lnot !A$}
\AxiomC{$\Gamma$}
\RightLabel{$\delta_1$}
\DeduceC{$!A$}
\RightLabel{\Elim{\lnot}}
\BinaryInfC{$\lfalse$}
\end{prooftree}
नव्या !!{derivation} मध्ये गृहीतक~$!A$ हे !!{discharged} आहे; म्हणून ही
$\Gamma$ पासूनची !!a{derivation} आहे.
\end{proof}

\begin{prop}
\ollabel{prop:prov-incons}
$\Gamma \Proves !A$ तेव्हा आणि केवळ तेव्हाच, जेव्हा $\Gamma \cup \{\lnot !A\}$
विसंगत असते.
\end{prop}

\begin{proof}
प्रथम $\Gamma \Proves !A$ असे समजा; म्हणजे $\Gamma$ मधील
!!{undischarged} गृहीतकांपासून $!A$ ची !!a{derivation}~$\delta_0$ असते.
पुढीलप्रमाणे $\Gamma \cup \{\lnot !A\}$ पासून $\lfalse$ ची
!!a{derivation} मिळते:
\begin{prooftree}
  \AxiomC{$\lnot !A$}
  \AxiomC{$\Gamma$}
  \RightLabel{$\delta_0$}
  \DeduceC{$!A$}
  \RightLabel{\Elim{\lnot}}
  \BinaryInfC{$\lfalse$}
\end{prooftree}

आता $\Gamma \cup \{\lnot !A\}$ विसंगत आहे असे समजा आणि $\Gamma \cup
\{\lnot !A\}$ मधील !!{undischarged} गृहीतकांपासून $\lfalse$ ची संबंधित
!!{derivation}~$\delta_1$ असू द्या. $\FalseCl$ वापरून पुढीलप्रमाणे केवळ
$\Gamma$ पासून $!A$ ची !!a{derivation} मिळते:
\begin{prooftree}
  \AxiomC{$\Gamma, \Discharge{\lnot !A}{1}$}
  \RightLabel{$\delta_1$}
  \DeduceC{$\lfalse$}
  \RightLabel{\FalseCl}
  \DischargeRule{\FalseCl}{1}
  \UnaryInfC{$!A$}
\end{prooftree}
\end{proof}

\begin{prob}
$\Gamma \Proves \lnot !A$ तेव्हा आणि केवळ तेव्हाच, जेव्हा $\Gamma \cup \{!A\}$
विसंगत असते, हे सिद्ध करा.
\end{prob}

\begin{prop}\ollabel{prop:explicit-inc}
  जर $\Gamma \Proves !A$ आणि $\lnot !A \in \Gamma$ असेल, तर $\Gamma$
  विसंगत आहे.
\end{prop}

\begin{proof}
  समजा $\Gamma \Proves !A$ आणि $\lnot !A \in \Gamma$. मग $\Gamma$ पासून
  $!A$ ची !!a{derivation}~$\delta$ असते. $\Elim{\lnot}$ नियमाचा पुढील सोपा
  उपयोग विचारात घ्या:
  \begin{prooftree}
    \AxiomC{$\lnot !A$}
    \AxiomC{$\Gamma$}
    \RightLabel{$\delta$}
    \DeduceC{$!A$}
    \RightLabel{\Elim{\lnot}}
    \BinaryInfC{$\lfalse$}
  \end{prooftree}
  $\lnot !A \in \Gamma$ असल्यामुळे सर्व !!{undischarged} गृहीतके
  $\Gamma$ मध्ये आहेत. म्हणून यावरून $\Gamma \Proves \lfalse$ हे दिसते.
\end{proof}

\begin{prop}\ollabel{prop:provability-exhaustive}
  $\Gamma \cup \{!A\}$ आणि $\Gamma \cup \{\lnot !A\}$ हे दोन्ही संच
  विसंगत असतील, तर $\Gamma$ विसंगत आहे.
\end{prop}

\begin{proof}
$\delta_1$ आणि $\delta_2$ या अनुक्रमे $\Gamma \cup \{ !A \}$ पासून
$\lfalse$ ची आणि $\Gamma \cup \{ \lnot !A \}$ पासून $\lfalse$ ची
!!{derivation}s असू द्या. मग आपण पुढीलप्रमाणे !!{derive} करू शकतो:
\begin{prooftree}
\AxiomC{$\Gamma, \Discharge{\lnot !A}{2}$}
\RightLabel{$\delta_2$}
\DeduceC{$\lfalse$}
\DischargeRule{\Intro{\lnot}}{2}
\UnaryInfC{$\lnot \lnot !A$}
\AxiomC{$\Gamma, \Discharge{!A}{1}$}
\RightLabel{$\delta_1$}
\DeduceC{$\lfalse$}
\DischargeRule{\Intro{\lnot}}{1}
\UnaryInfC{$\lnot !A$}
\RightLabel{\Elim{\lnot}}
\BinaryInfC{$\lfalse$}
\end{prooftree}
$!A$ आणि $\lnot !A$ ही गृहीतके !!{discharged} असल्यामुळे ही केवळ
$\Gamma$ पासून $\lfalse$ ची !!a{derivation} आहे. म्हणून $\Gamma$ विसंगत
आहे.
\end{proof}

\end{document}
